\documentclass{report}
\usepackage{amsmath,amssymb}
\setlength{\parindent}{0mm}
\setlength{\parskip}{1em}
\begin{document}
\begin{center}
\rule{6in}{1pt} \
{\large Stefan M Wild \\
{\bf Manifold Sampling for L1 Nonconvex Optimization}}

Argonne National Laboratory \\ Mathematics and Computer Science Division \\ 9700 S Cass Ave \\ Bldg 240-1151 \\ Lemont \\ IL 60439
\\
{\tt wild@anl.gov}\\
Jeffrey Larson\\
Matt Menickelly\end{center}

We introduce the \emph{manifold sampling} algorithm for
unconstrained minimization of
a nonsmooth composite function $h\circ F$. By classifying points in the domain
of $h$ into manifolds, the algorithm adapts search directions within a
trust-region framework based
on knowledge of manifolds
intersecting the current trust region.

We motivate this idea through a study of $\ell_1$ functions, where the
classification into manifolds using zero-order information about the
constituent functions $F_i$ is trivial.
In this case, we show that both deterministic and stochastic versions of the
algorithm generate cluster points that are Clarke stationary.


Since the manifold sampling algorithm only requires approximations of the
Jacobian $\nabla F$ at the trust-region center, the algorithm can also address
situations where $F$ is defined by a simulation for which such derivative
information may not be available.
We present numerical results for different variants of the algorithm and
show that
using manifold information from points near the current iterate can improve
empirical performance.


\end{document}
